% ID: MAT-TRI-EQG-015
% Titolo: Equazioni goniometriche con fattorizzazione (1)
% Disciplina: Matematica
% Area: Trigonometria
% Argomento: Equazioni goniometriche
% Classe: Quarta secondaria di secondo grado
% Difficolta: 3
% Tipo: problema_numerico
% Tag: matematica, trigonometria, equazioni_goniometriche
% Fonte: verifica_goniometria_base_anonimizzata
% Provenienza: materiale fornito dal docente; testo anonimizzato; soluzione redatta per Exergo

\begin{esercizio}
Risolvi le seguenti equazioni: 

\begin{enumerate}
\item $\displaystyle\sin 2x=\sin x;$

\item $\displaystyle 2\cos^2 x+3\cos x+1=0.$
\end{enumerate}
Salvo diversa indicazione, lavora in $\mathbb R$; angoli in radianti.
\end{esercizio}

\begin{soluzione}
a) $\sin x(2\cos x-1)=0$:
$x=k\pi$ oppure $x=\pm\pi/3+2k\pi$.
b) $(2\cos x+1)(\cos x+1)=0$:
$x=2\pi/3+2k\pi,4\pi/3+2k\pi,\pi+2k\pi$. In tutte le famiglie di soluzioni $k\in\mathbb{Z}$.
\end{soluzione}
